% 本文件由 LaTeX 排版 Agent 根据有效 translation.json 自动生成。
% author=Mar Jacob; work_id=0860; title=殉道者哈比卜颂

\chapter{殉道者哈比卜颂}

% structure: p0001 source=h1 target=omitted reason=page-title-already-rendered-by-parent

\Emph{\WorkTitle{殉道者哈比卜颂}，\Person{雅各伯}所作。}\par

殉道者\Person{哈比卜}，身披火焰，从火中向我呼唤，\par

要我为他也在荣耀者中塑造一幅美丽的肖像。\par

征服者的同伴，看！他从烈焰中向我招手，\par

要我因他主的荣耀，歌唱赞美他。\par

在燃烧的炭火中，那位\Emph{勇士}站立，看！他向我呼唤，\par

要我塑造他的肖像：但熊熊烈火却不允许。\par

他的爱炽热，他的信德发光；\par

他的火也在燃烧，谁能充分述说他的爱？\par

不，因着那引领殉道者进入烈火的爱，\par

无人能够述说他那神圣的美。\par

因为谁敢进入那烈火中观看，\par

他像谁，在荣耀者中应按何种样式塑造？\par

我是否应将他的肖像塑造在那些青年，即火窑之子旁边？\par

我应将\Person{哈比卜}与\Person{阿纳尼雅}并列吗？我不知道。\par

看！他们并未在那里被焚烧：那么，他怎能相似他们？\par

\Emph{我说}，他像他们，当他被焚烧而青年们却未被焚烧时？\par

我问，谁\Emph{更美丽}——殉道者\Person{哈比卜}，还是\Person{阿匝黎雅}？\par

这肖像对我而言是难的：我该如何观看，我不知道。\par

看！\Person{米加耳}未被火焰焚烧；\par

但\Person{哈比卜}被焚烧了：那么，在观看者眼中，谁\Emph{更美丽}？\par

谁敢说这个可厌，或那个可厌；\par

或在观看者眼中，这个不如那个俊美？\par

\Emph{有三个}在火中，火焰不接近他们；\par

但有一个被焚烧了：我怎能足够地讲述\par

那第四位的\Emph{形状}就是那下到火窑中的那一位的，\par

为要与\Emph{那三个}一同为\Person{哈比卜}塑造一幅肖像？\par

祂在火中给被焚烧者一个位置，\par

使他能代替那位第四位，与征服者并肩。\par

如果那三个虽未被焚烧，其荣美已是荣耀的，\par

那么这位\Emph{被焚烧}的，为何不能与荣耀者相融？\par

若一个人有能力选择被焚烧或不焚烧，\par

那么这位\Emph{被焚烧}的，其美比那三个更为高超。\par

但既然主是\Emph{万有}的主宰，\par

无论祂拯救还是交付，祂都当受赞美。\par

此外，那三个未被焚烧者的意愿，\par

与那位被焚烧者的意愿，在此处和在彼处，都是同一的。\par

若火的主命令火焚烧他们，\par

那三个\Emph{也}必被焚烧；\par

若祂示意火不要焚烧那一个人，\par

他也不会被焚烧；他得救也并非出于自己。\par

进入火中是出于他们自己的意愿，当他们进入时；\par

但他们未被焚烧——\Emph{因为}火的主愿意并命令了。\par

因此，被焚烧者的美与未被焚烧者的美是相等的，\par

因为意愿也是相等的。\par

可爱的殉道者\Person{哈比卜}！你的美是崇高的，你的品位是崇高的：\par

你的冠冕也是优雅的，你的事迹与荣耀者并列。\par

你是精选的金子，火已试炼了你，你的荣美灿烂。\par

看！你被编入君王的冠冕中，与得胜者一同。\par

在\Person{天主圣子}的教导中，\par

像勇士般前行，因他信德之美。\par

殉道者\Person{哈比卜}是一位真理的导师；\par

也是一位宣讲者，他的口充满信德。\par

他警醒且\Emph{敏捷于服务}，以他的教导鼓励\par

天主的家室，藉着他的信德。\par

他充满光明，与黑暗搏斗，\par

这黑暗从异教而来，笼罩了大地。\par

在聚会中，他的口充满圣子的福音；\par

当他到达村庄时，他几乎成了领路者。\par

他热心，因为他关心神圣的教导，\par

为坚固信德的追随者。\par

当异教的风吹起时，他是一盏灯，\par

当他们吹向他时，他发出火焰，并未熄灭。\par

他全身燃烧，充满对他主的爱，并关心\par

这一点——他能毫无阻碍地谈论祂。\par

错误的荆棘从异教中在地上长出；\par

他尽己所能，以勤奋拔除它们。\par

他教导、劝诫并坚固信德，\par

基督的朋友们，受逼迫者所困扰。\par

他既与刀剑也与火焰搏斗，\par

以炽热如焰的爱，毫不畏惧。\par

像一把双刃的剑，锋利的是\par

他的信德，他与错误争战。\par

他成为这地的酵母，这地已因\par

对虚空偶像的迷恋而疲乏，这迷恋是错误带来的。\par

他像盐，以他可口的教导\par

供给这地区，这地区因不信而变得乏味。\par

他是一位执事，却履行了大司祭的职责，\par

藉着宣讲和教导真理。\par

当他管理羊群时，他是羊群的好牧人；\par

他牧养时，为羊群舍命。\par

他赶走豺狼，驱离野兽。\par

他修补破口，将羔羊聚回羊圈。\par

他秘密出去鼓励会众：\par

他坚固他们，劝勉他们，扶持他们。\par

他锻造信德的盔甲，给他们穿上，\par

免得他们被盛行的异教可耻地击败。\par

天主圣子羊圈的羊群正被\par

逼迫者蹂躏：他鼓励羊羔和母羊。\par

他是信德家室的辩护者；\par

他教导他们不要被逼迫者吓倒。\par

他教导他们奔迎死亡，\par

毫不畏惧刀剑或火焰。\par

在天主圣子的教导中他兴盛，\par

以致他的信德无畏前行。\par

然后错误嫉妒，因他而暴怒、疯狂；\par

她追赶他，要在世上流无辜的血。\par

那憎恨人类的毁谤者，\par

为他设下陷阱，要除掉他的存在。\par

那憎恨真理的追赶他，要置他于死，\par

为使他停止在天主家室中的教导之声。\par

错误呼喊，\Emph{要求}\Person{哈比卜}死，因她恨他；\par

烦恼催逼着她，她想要取他的性命。\par

他的事迹传到了当地外教法官的耳中，\par

他珍贵的美名也传到了国王那里：国王盛怒之下，\par

因王冠与外教交织在一起，便下令处死\par

哈比卜，因为他充满了信德。\par

命令下达法官时，他便武装自己，\par

满腔怒火与暴烈；带着嗜血的心，\par

如同猎人设网捕捉幼鹿，\par

他们出去搜捕哈比卜，要抓住\Emph{他}。\par

但此人是一位信德的宣讲者，\par

在十字架的大道上兴旺发展；\par

为借教导裨益自己民族的子民，\par

他的工作遍及周围各地。\par

因此，当谬误出去追捕他时，竟寻他不着：\par

并非因他逃跑，而是因他出去宣讲福音了。\par

那时，因外教人的愤怒极其不当，\par

他们便为他的缘故抓住了他的亲属和母亲。\par

女子啊，你有福了！你既是殉道者的母亲。\par

他们为何不义地抓住\Emph{你}并捆绑你？\par

他们向你要求什么，你这充满美丽的人？\Emph{我问}，他们向你要求了什么？\par

看！他们要求你带来殉道者，使他成为祭献。\par

带来，带来你甘甜的果实到奉献的地方——\par

\Emph{这果实}香气芬芳，可作为向天主性的馨香。\par

秀美的枝条，从原处带来你的簇串，\par

使它的酒可作滋味甜美的奠祭。\par

羔羊听说他们正在寻找他，要献他为祭；\par

他便起身，欢欢喜喜地来到献祭者面前。\par

他听说别人也因他而受苦，\par

便前来代替许多人承受属他的苦难。\par

签落在他身上，要独自一人成为祭献；\par

那要献上他的烈火正等着\Emph{他}，直到他到来。\par

那些为他的缘故而被捆绑的许多人中，\par

没有一个人被捉去受死，惟独他一人。\par

他才是配得的，为他预备了殉道；\par

无人能夺去殉道者的位置。\par

因此，他甘愿将自己呈献\par

给法官，为被捉拿，并为耶稣的缘故而死。\par

他听说他们搜寻他，便前来被捉，正如他们所寻：\par

他自动来到法官面前，面容毫无畏惧。\par

他没有躲藏，也不愿逃离法官：\par

因为他浸透了光明，不愿逃离黑暗。\par

他不是\Emph{强盗}，不是凶手，不是窃贼，\par

不是黑夜之子：他的一切行程都在光天化日之下。\par

好牧人为何要逃离他的羊群，\par

让羊栈被强盗吞噬？\par

医生为何要逃跑，他本出去医治疾病，\par

并用天主子之血治愈灵魂？\par

这\Emph{勇敢}之人带着无畏的面容和伟大的心；\par

他欢欢喜喜地跑去迎接死亡，为耶稣的缘故。\par

他进去，站在法官面前，对他说：\par

我就是你搜寻的哈比卜：看！\Emph{我}站在这儿。\par

外教人颤抖，惊愕抓住他，他对他感到惊奇——\par

对那既不畏惧刀剑也不畏惧烈火的人。\par

正当他以为他正在急速逃跑时，他却进来嘲弄了他；\par

法官战栗，因为见他即使在\Emph{死亡}面前也如此勇敢。\par

他是那位天主子——祂曾说——的门徒：\par

\Quote{起来，我们走吧。看，那出卖我的人近了。}\par

祂又对钉死祂的人说：\Quote{你们找谁？}\par

他们说：\Quote{耶稣。}祂对他们说：\Quote{我就是。}\par

天主子甘愿来到十字架上；\par

殉道者注视着祂，便\Emph{毫不勉强}地把自己呈现在法官面前。\par

外教人看见他，便恐惧战栗，\Emph{对他}恼怒。\par

他的怒气被激发，便在狂怒中开始向他提问。\par

仿佛他曾将被杀者的血洒在地上，\par

他继续审问这位圣者，但圣者并不羞愧：\par

一面威胁他，设法恐吓他，惊吓他，\par

一面数算他为他预备的苦难。\par

但哈比卜受审时不害怕，\par

不羞愧，也不被所\Emph{听到}的威吓吓倒。\par

他扬声承认耶稣是天主子——\par

说自己是祂的仆人，是祂的司铎，是祂的侍奉者。\par

外教人狂怒，像狮子般向他咆哮，\par

他不颤栗，也不停止承认天主子。\par

他受鞭打，鞭伤对他极为宝贵，\par

因他承担了天主子所受的一点点创伤。\par

他被捆锁，他仰望他的主——他们也捆绑了祂；\par

他的心欢喜，因他开始行走在祂苦难的道路上。\par

他上了刑架，他们用铁梳撕裂他，但他的灵魂闪烁着光明，\par

因为他被\Emph{认为}配得上让十字架苦难的痛苦临到他身上。\par

他已决心行走在死亡的道路上，\par

他在其中还能期望找到什么，除了苦难？\par

祭献的烈火已与他订婚，他期待着她；\par

而烈火\Emph{方面}则遣来铁梳、鞭伤和痛苦，让他尝受。\par

在她来临的整个期间，她差遣苦难给他，为借这些苦难\par

使他预备好，好使当她遇见他时，不致让他惊慌。\par

苦难炼净了他，好使当熊熊烈火考验他时，\par

在他精选的黄金中\Emph{找不到}任何渣滓。\par

他忍受了临到他的一切痛苦，\par

为体验\Emph{苦难}，并在火焰中像勇士般站立。\par

他欢欢喜喜地接受了必须承受的苦难：\par

因他知道，在它们结束后，他将找到死亡。\par

他不惧怕死亡，也不惧怕苦难：\par

因他的心陶醉于十字架之酒。\par

他轻视他的身体，任其被迫害者拖曳；\par

也轻视他的四肢，在剧痛中被撕裂。\par

背上的鞭打，肋旁的铁梳，脚上的枷锁，\par

面前的烈火：他仍然勇敢且充满信德。\par

他们嘲讽他：看！你敬拜一个人；\par

但他却说：我不敬拜人，\par

而是敬拜那取了身体并成为人的天主：\par

我敬拜祂，因为祂与生祂者同为天主。\par

殉道者哈比卜的信德充满光明\par

有信德之\Emph{城}厄德撒被光照。\par

阿布加尔之女，阿戴将之许配给十字架——\par

赖她有光明，赖她有真理与信德。\par

她的君王来自此，她的殉道者来自此，她的真理来自此；\par

\Emph{她的}信德的教师们也来自此。\par

阿布加尔相信你是天主，天主之子；\par

因他信德之美，他获得了祝福。\par

厄德撒人之子、殉道者沙尔比勒更曾说：\par

我的心被降生成人的天主所俘获。\par

而殉道者哈比卜，也曾在厄德撒受冠冕，\par

他宣认了这些：祂取了身体，降生成人；\par

祂是天主之子，也是天主，并且降生成人。\par

厄德撒从教师那里学到了真实的事：\par

她的君王教导她，她的殉道者教导她信德；\par

但对其他那些欺诈的教师，她不肯听从。\par

殉道者哈比卜，在厄德撒耳中如此高声呼喊，\par

从火焰中：我不敬拜人，\par

而是天主，祂取了身体并降生成人，\par

我敬拜祂。殉道者扬声\Emph{如此}宣认。\par

从被铁梳撕裂、被焚烧、\Emph{被抬上刑台}、被杀害的宣信者，\par

并\Emph{从}一位正义的君王，厄德撒学到了信德，\par

她认识我们的主——祂就是天主，天主之子；\par

她也学到了并坚信祂取了身体并降生成人。\par

她不是从普通的书吏学到信德：\par

她的君王教导她，她的殉道者教导她；她坚信他们：\par

若有人毁谤她曾敬拜过一个人，\par

她便指向她的殉道者，他们为祂而死，因为祂是天主。\par

我不敬拜人，哈比卜说，\par

因为经上记载：\Quote{凡信赖人的，是可咒骂的。}\par

既然祂是天主，我敬拜祂，是的，情愿被烧死\par

为祂的缘故，决不背弃祂的信德。\par

这真理，厄德撒自幼就持守，\par

在她年老时，她不会像穷人之女那样将它卖掉。\par

她的正义君王成了她的书吏，从他她学到了\par

关于我们的主——祂是天主之子，也是天主。\par

阿戴，带来了新郎的戒指，戴在她手上，\par

就这样将她许配给天主之子，那\Emph{独生}子。\par

司祭沙尔比勒，曾试验并验证了所有神祇，\par

他死了，正如他所说，\Quote{为那降生成人的天主。}\par

沙穆纳和古里亚，为了\Emph{独生}子的缘故，\par

伸颈\Emph{受斩}，为祂而死，因为祂是天主。\par

而殉道者哈比卜，曾是会众的教师，\par

宣讲祂，说祂取了身体并降生成人。\par

为一个人，殉道者不会\Emph{甘愿}被火烧死；\par

但他被烧死，是\Quote{为了那降生成人的天主。}\par

厄德撒是见证，他受火刑时如此宣认：\par

从被烧死的殉道者的宣认，谁能逃脱呢？\par

信德使所有理智缄默并轻蔑——\par

她充满光明，不屈从于阴影。\par

她轻蔑那否认圣子是天主而毁谤祂的人；\par

也轻蔑那说\Quote{祂未取身体，也未降生成人}的人。\par

在充满真理的信德中，他站在火上；\par

他成了乳香，以他的馨香平息了天主之子。\par

在他\Emph{的}所有痛苦中，在他\Emph{的}所有折磨中，在他\Emph{的}所有苦难中，\par

他如此宣认，也如此教导那蒙福之城。\par

厄德撒紧紧持守这关于我们主的真理——\par

即祂是天主，由玛利亚降生成人。\par

新娘恨那否认祂天主性的人，\par

她轻视并藐视那毁谤祂有形身体的人。\par

她承认祂在天主性和人性中是\Emph{一体}——\par

那\Emph{独生}子，祂的身体与祂不可分离。\par

帕提亚的女儿就这样学会了信仰，\par

她就这样坚定持守，也这样教导那听从她的人。\par

因此，充满\Emph{外教热忱}的审判官下令，\par

将殉道者带出去，在为他预备的火中焚烧。\par

随即，一条皮带塞入他口中，如同\Emph{对待凶手}，\par

他的宣认藏于内心，向着天主。\par

他们催促他，他走出审判厅，欢欣踊跃，\par

因为那时刻已到，冠冕将赐予他的信德。\par

成群的人与他同去，为他作伴，\par

看着他，不是像一个死人\Emph{被送去埋葬}，\par

而是一个即将藉火焰成为新郎的人，\par

那按公义为他存留的冠冕即将赐下。\par

他们看他如同一个上战场的人，\par

周围是长矛、枪和剑，但他战胜了它们。\par

他们看见他像得胜的斗士一样上去，\par

在他胜利时，围观者向他献上花环。\par

他们看他如何战胜了掌权者和异能者，\par

它们都与他争战，而他使它们蒙羞。\par

基督追随者的整个会众因他而欢腾，\par

因为他所承受的苦难激励了信德的朋友。\par

教会与他同去，一位充满光明的新娘；\par

她的容光照耀着那与她结合的蒙爱殉道者。\par

于是他的母亲，因这是她儿子的婚宴，\par

穿上比平时更高贵的衣裳。\par

因简陋的衣衫不适合宴席大厅，\par

她穿上\Emph{华服}，全白，恰到好处。\par

爱来到战场争战——\par

在母亲灵魂中——\Emph{天性之爱}和\Emph{天主之爱}。\par

她看着儿子被带出去投入火焰；\par

因她心中有主的爱，她不受苦。\par

她母腹的牵念为它的果实呼喊；\par

但信德使它们沉默，喧嚣止息。\par

天性为那从她身上分离的肢体尖叫；\par

但主的爱陶醉了灵魂，使她不觉痛苦。\par

天性爱，但主的爱在争战中得胜，\par

在母亲灵魂内，使她不为心爱者悲伤。\par

她心中充满各种喜乐之情，代替了痛苦；\par

她身穿华丽衣裳，代替了哀悼。\par

她陪他出去受火刑，欣喜若狂，\par

因为对主之爱战胜了本性之爱。\par

她\Emph{穿着}白色，如同新郎，摆设了婚宴——\par

她，殉道者的母亲，因他而欢欣。\par

我们可以称这有福之人为\Quote{第二\Person{沙慕纳}}吗？\par

因为，即使七人被焚烧代替一人，她也心满意足。\par

她有一个，就把他交给火作食物；\par

而且，就像那一个，如果她有七个，她也会将\Emph{他们全部}交出。\par

他被投入火中，火焰在他周围燃起；\par

他的母亲望着，并不因他被焚烧而悲伤。\par

另一只眼，注视着看不见的事物，\par

在她的灵魂里，因此当他在被焚烧时她欢欣踊跃。\par

她注视殉道者冠冕上的光明宝石，\par

以及为他们苦难之后所存留的荣耀；\par

以及\Emph{注目于}他们通过苦难继承的应许福分，\par

以及\Emph{注目于}以光明装扮他们肢体的天主圣子，\par

以及\Emph{注目于}那永不朽坏之国度的各种美丽，\par

以及\Emph{注目于}为他们敞开以进入天主面前的宽阔之门。\par

殉道者的母亲在他被焚烧时就注视着这些，\par

她就欢喜，她就昂首，并穿着白色与他一同前行。\par

她注视着他，当火焰吞噬他的身体时，\par

并且，因为他的冠冕非常高贵，她并不悲伤。\par

甘甜的根被投入火中，落在炭火上；\par

它变成了乳香，洁净了空气中的污秽。\par

空气曾被祭献的烟雾污染，\par

而藉着这位殉道者的被焚，空气得以洁净。\par

穹苍因祭台的烟气而恶臭；\par

殉道者的馨香升起，空气因而变得甘甜。\par

祭献停止了，集会中有了和平；\par

刀剑变钝了，不再蹂躏基督的朋友们。\par

在我们的土地上，从\Person{沙比尔}开始，以\Person{哈比卜}结束；\par

自从他被焚以来，直到现在，它没有杀死一个人。\par

\Person{君士坦丁}，征服者之首，取得了帝国，\par

十字架践踏了皇帝的冠冕，并安放在他头上。\par

偶像崇拜的高傲之角被折断，\par

自从殉道者被焚以来直到现在，它没有刺穿一人。\par

他的烟雾升起，成了对天主性的馨香；\par

藉此，被异教污染的空气得到净化，\par

并藉着他的被焚，整个土地被洁净：\par

愿那赐予他冠冕、荣耀和美名者受赞颂！\par

\Emph{于此}结束由\Person{玛尔雅各布}所作的\WorkTitle{关于殉道者哈比卜的讲道}。\par

\SourceRecord{Translated by B.P. Pratten.；Ante-Nicene Fathers；Vol. 8.；Edited by Alexander Roberts, James Donaldson, and A. Cleveland Coxe.；Buffalo, NY: Christian Literature Publishing Co.,；1886.；<http://www.newadvent.org/fathers/0860.htm>.}
